% NOTE: unlike Sections 2-6, this file has NO source docs/SECTION_1_DRAFT.md
% to port from -- the task that names it ("the human has committed ...
% sections/1_introduction.tex") assumed a file that was never actually
% pushed to the repository. Drafted here from docs/PAPER_SKELETON.md's
% Section 1 outline instead, flagged for human review since it is new
% prose, not a port.
\section{Introduction}
\label{sec:introduction}

Emotion recognition in conversation (ERC) underlies a wide range of applications, from moderating and summarizing customer-support interactions to flagging distress in mental-health and crisis-line transcripts to giving conversational agents a sense of a user's affective state. Because these are conversations rather than isolated utterances, the field's architectural trajectory has consistently added conversational structure on top of per-utterance representations: recurrent per-speaker state, then graph formulations that treat utterances or speakers as nodes and propagate influence between them, including DialogueGCN~\cite{dialoguegcn2019}, MMGCN~\cite{mmgcn2021}, COGMEN~\cite{cogmen2022}, and our own prior SocialArcNet~\cite{socialarcnet2026} (published: $0.62$ weighted F1 on MELD). This machinery has consistently earned its keep in the literature that reports it.

That literature, however, shares an experimental contract that has gone largely unexamined: the graph or recurrent architecture is almost always trained over frozen, or only lightly adapted, utterance encoder features. This was a pragmatic necessity when fine-tuning large encoders was expensive; it no longer is. Parameter-efficient adaptation methods make it routine to fine-tune a strong text encoder on a single consumer GPU in minutes per epoch. This raises the question this paper answers directly: does conversation-graph machinery retain measurable value once the utterance encoder it sits on top of is itself fine-tuned and conversationally aware, or does the machinery's reported value substitute for a representational gap that fine-tuning now closes on its own?

Answering this cleanly requires an attribution instrument that ordinary ablation studies do not provide. Retraining an architecture from scratch with and without a component conflates the component's true contribution with the incidental cost of its presence: extra randomly initialized parameters that must be optimized before the model even recovers baseline competence, altered optimization dynamics, and additional overfitting capacity on corpora that are not large. We show empirically (Section~\ref{sec:meld-mechanism}) that this confound is not a theoretical nicety: it inflates the measured harm of adding components by over $60\%$ relative to a construction that removes it. This paper makes four contributions:

\begin{itemize}
\item[\textbf{C1}] A residual attribution methodology in which every architectural component provably initializes as an identity function on a frozen, fine-tuned foundation, so that any measured delta is attributable to the component by construction rather than by ceteris-paribus hope (Section~\ref{sec:methodology}).
\item[\textbf{C2}] The fine-tuning boundary finding: conversation-graph machinery's gains, real and reproducible on frozen features, vanish (and turn slightly negative) once the text encoder is fine-tuned, even when that encoder is given no conversational context at all; a context-window sweep further shows that fine-tuning itself, not the size of the context window, is the operative variable (Section~\ref{sec:meld}).
\item[\textbf{C3}] A fair, pre-registered trial of relational (pair-state) memory and a shift-aware auxiliary objective on IEMOCAP's long dyadic dialogues (the corpus, protocol, and label chosen in advance specifically to favor the relational mechanism), whose result refutes rather than rescues it, and whose mechanism analysis identifies a genuine dissociation between what the component learns and what the classification task can use (Section~\ref{sec:iemocap}, Section~\ref{sec:analysis}).
\item[\textbf{C4}] A fully pre-registered, seed-replicated, publicly versioned experimental protocol (gates, locked recipes, and a per-decision audit trail) released alongside the paper as a reusable template for attribution-honest architecture evaluation (Section~\ref{sec:protocol}).
\end{itemize}

Taken together, these results do not indict the graph-ERC literature's original findings, which were real on the foundations of their day; they indict a default evaluation practice that has not kept pace with how cheap strong foundations have become. Our closing counsel to the field, developed across Sections~\ref{sec:meld}--\ref{sec:analysis} and restated in Section~\ref{sec:conclusion}, is a change of default: architectural contributions to ERC should be attributed against a competitively fine-tuned encoder baseline, using an instrument that separates a component's true contribution from the cost of its own re-training, before a frozen-feature comparison is treated as sufficient evidence of the component's value.
